\chapter*{Lời Mở Đầu}
\addcontentsline{toc}{chapter}{Lời Mở Đầu}
\markboth{Lời Mở Đầu}{Lời Mở Đầu}

\begin{truyen}{Lời Mở Đầu}{Opening Words}
\chuhoa{Đ}{ạo làm người không học một lần rồi thôi — nó cần được nhắc đi nhắc lại suốt đời.}
\textit{(The way of being human is not learned once and forgotten --- it needs to be revisited throughout life.)}

Quyển hai của bộ sách này tiếp tục hành trình đó. Những câu chuyện ở đây xoay quanh đạo làm con, làm bạn, làm trò, làm thầy, làm người.
\textit{(The second volume of this series continues that journey. The stories here revolve around being a good child, friend, student, teacher, and human being.)}

Mỗi bài học được đặt trong một tình huống quen thuộc để người đọc dễ liên hệ với cuộc sống thực của chính mình.
\textit{(Each lesson is placed within a familiar situation so readers can easily relate it to their own real lives.)}

Đọc sách này như nghe lại những lời căn dặn ta đã biết nhưng hay quên.
\textit{(Read this book as if hearing again the reminders you already knew but often forget.)}
\end{truyen}

\begin{baihoc}
Đạo làm người không khó --- nhưng cần được nhắc nhở và thực hành mỗi ngày.
\textit{(The way of being human is not difficult --- but it needs daily reminders and practice.)}
\end{baihoc}

\begin{ghinhoanh}
\textbf{Short English to remember:}

Character is built one story at a time.

Being good takes daily practice, not one effort.
\end{ghinhoanh}

\begin{tuhocnhanh}
1. Trong các đạo lý ở quyển này, đạo nào em thấy mình còn thiếu nhất? \textit{Among the values in this volume, which do you feel you still lack the most?}

2. Ai trong cuộc sống em là hình mẫu cho một trong những đạo này? \textit{Who in your life is a model for one of these values?}

3. Em có thể làm gì hôm nay để sống đúng với đạo mình đang học? \textit{What can you do today to live in alignment with the value you are studying?}
\end{tuhocnhanh}

\ngancach
